% !TeX root = ../tesi.tex

\chapter*{Sommario}

In questa relazione descrivo lo \emph{stage} che ho svolto presso Halue S.r.l., società di
consulenza tecnologica con sede operativa ad Abano Terme, dal 14 settembre al 9 novembre 2026.

Con la supervisione del Prof. Tullio Vardanega, relatore, e di Antonio Cavestro, \emph{tutor}
aziendale, ho progettato e realizzato ChatCart, il prototipo di un assistente conversazionale per i
negozi Shopify che risponde al cliente non soltanto con testo, ma anche con componenti
dell’interfaccia che sceglie e popola in tempo reale con i dati del negozio.

\textbf{Struttura del documento.} La relazione si articola in quattro capitoli:

\begin{itemize}
        \item \textbf{Capitolo 1: «Halue e il commercio digitale».} Presento l’azienda, la sua
	clientela, le tecnologie e i processi che adotta e il suo rapporto con l’innovazione, in
	particolare con l’intelligenza artificiale.
        \item \textbf{Capitolo 2: «Oltre la conversazione testuale».} Discuto il problema che lo
	\emph{stage} affronta, il suo posto nella strategia di Halue, gli obiettivi e i vincoli del
	progetto e le ragioni per cui l’ho scelto.
        \item \textbf{Capitolo 3: «Progetto e sviluppo di ChatCart».} Illustro il metodo di lavoro,
	l’analisi, l’architettura, la realizzazione dell’agente e del \emph{widget}, le funzionalità
	sperimentali e le possibili evoluzioni del progetto, la verifica e i risultati ottenuti.
        \item \textbf{Capitolo 4: «Valutazione retrospettiva».} Valuto il grado di soddisfacimento
	degli obiettivi, la mia maturazione professionale e la distanza tra le competenze che lo
	\emph{stage} mi ha richiesto e quelle che ho acquisito nel corso di studi.
\end{itemize}

In appendice raccolgo gli acronimi e le abbreviazioni con la loro forma estesa, il glossario e i riferimenti.

\newpage

\textbf{Convenzioni tipografiche.} Nel documento adotto le seguenti convenzioni:

\begin{itemize}
        \item il corsivo indica i termini in lingua straniera, e soltanto questi; i nomi propri di
	aziende, prodotti e metodi (come Shopify, Slack o Scrum) restano in tondo;
        \item alla prima occorrenza di un termine presente nel glossario ne do nel testo una
	definizione essenziale e lo contrassegno con il pedice G; il termine e il pedice sono cliccabili e
	rimandano all’approfondimento nel glossario;
        \item figure, tabelle e codici sorgente hanno una numerazione progressiva e una didascalia
	posta sotto l’elemento, che ne indica la fonte quando non sono di mia realizzazione.
\end{itemize}
